\chapter{Mix-FFN, LayerNorm, residuals and the transformer block}\label{ch:block}

\begin{onething}
A SegFormer block is two sublayers, attention and Mix-FFN, each wrapped the same way: normalise
the input with LayerNorm, transform it, drop a random part of the result during training, and add
it back to what came in. Two lines of code, and every piece has a reason you can calculate.
\end{onething}

\Cref{ch:attention,ch:esa} built the attention half of the block. This chapter builds the other
half, the Mix-FFN, and the three pieces of wiring that hold the two together: LayerNorm, the
residual connection and dropout. It ends by counting the parameters of a whole block by hand and
matching the real model's number exactly. The code is
\file{app/foundation_models/components/mix_ffn.py} and
\file{app/foundation_models/components/segformer_block.py}; the numbers are printed by
\file{research/textbook/tools/compute_13-15.py}.

\section{Why attention alone is not enough}\label{sec:block:why}

Three problems remain after \cref{ch:esa}.

\paragraph{1. Attention gathers but hardly computes.} Look at what attention does to the values:
it multiplies them by the weights matrix $\mat{A}$ and then by $\mat{W}_O$. If the weights were
held fixed, the output $\mat{A}\mat{X}\mat{W}_V\mat{W}_O$ would be a \emph{linear} function of the
input: doubling every token would double every output. The only nonlinearity is the softmax, and
it acts on the weights, deciding \emph{who} is listened to, not on the content that is passed
along. A stack of linear maps is still a linear map ($\mat{W}_1\mat{W}_2$ is one matrix), and a
linear map cannot represent ``bright in band 40 \emph{and} dark in band 90''. Each token needs some
private, nonlinear processing. That is the feed-forward network.

\paragraph{2. Deep stacks are hard to train.} The encoder has eight blocks, each with two
sublayers. Training sends an error signal backwards through all of them, and every sublayer
multiplies it by its own derivative. Sixteen multiplications can shrink the signal to nothing or
blow it up. The residual connection fixes this.

\paragraph{3. Scales drift.} Each sublayer adds its output to the stream of tokens, so the typical
size of a token's numbers changes from block to block. \Cref{sec:attention:scale} chose the divisor
$\sqrt d$ assuming unit-variance inputs; if the inputs drift, the softmax saturates after all.
LayerNorm puts every token back on a standard scale before each sublayer sees it.

\begin{analogy}
Think of a document passed around an editorial office. Each editor receives a clean photocopy
(LayerNorm: everyone works from a standardised copy, whatever state the original is in), writes
suggested changes in the margin (the sublayer), and some suggestions are struck out at random
during training so that nobody relies on one editor (dropout). The suggestions are then applied to
the \emph{original} (the residual addition): the document is never replaced by an editor's
rewrite, only amended. Even if one editor does nothing useful, the document reaches the next desk
intact.
\end{analogy}

\section{LayerNorm}\label{sec:block:ln}

\subsection{The derivation}

LayerNorm\index{LayerNorm} \citep{ba2016layer} standardises each token on its own, across its $C$
features. Take one token $\vect{x} = (x_1, \dots, x_C)$.

\emph{Step 1, the mean} of its features:
\begin{equation}\label{eq:block:mu}
  \mu = \frac{1}{C}\sum_{i=1}^{C} x_i .
\end{equation}
\emph{Step 2, the variance}, the average squared distance from that mean (divided by $C$, not
$C-1$: PyTorch's LayerNorm uses this ``biased'' form):
\begin{equation}\label{eq:block:var}
  \sigma^2 = \frac{1}{C}\sum_{i=1}^{C} (x_i - \mu)^2 .
\end{equation}
\emph{Step 3, normalise:} subtract the mean, so the features centre on zero, and divide by the
standard deviation, so their spread becomes one. A small $\epsilon$ (PyTorch's default is
$10^{-5}$) prevents division by zero when all features are equal:
\begin{equation}\label{eq:block:norm}
  \hat x_i = \frac{x_i - \mu}{\sqrt{\sigma^2 + \epsilon}} .
\end{equation}
\emph{Step 4, the learned affine map.} Forcing every token to mean 0 and spread 1 may throw away
something useful, so each feature gets a learned scale $\gamma_i$ and shift $\beta_i$:
\begin{equation}\label{eq:block:ln}
  \LN(\vect{x})_i = \gamma_i\,\hat x_i + \beta_i .
\end{equation}
The vectors $\boldsymbol{\gamma}$ and $\boldsymbol{\beta}$ are the only parameters, $C$ each, so
\begin{equation}\label{eq:block:lnparams}
  \text{LayerNorm}(C) \text{ has } 2C \text{ parameters}.
\end{equation}
They start at $\gamma_i = 1$, $\beta_i = 0$, which makes LayerNorm a pure standardisation at
initialisation.

\begin{worked}[label={we:block:ln}]{LayerNorm of a 4-number token}
Take $\vect{x} = (1.0, 2.0, 3.0, 4.0)$ with $\boldsymbol\gamma = (1,1,1,1)$ and
$\boldsymbol\beta = (0,0,0,0)$. Do it on paper.

\emph{Mean.}
\[
  \mu = \frac{1.0 + 2.0 + 3.0 + 4.0}{4} = \frac{10.0}{4} = 2.5 .
\]
\emph{Deviations and their squares.}
\begin{align*}
  1.0 - 2.5 &= -1.5, & (-1.5)^2 &= 2.25\\
  2.0 - 2.5 &= -0.5, & (-0.5)^2 &= 0.25\\
  3.0 - 2.5 &= \phantom{-}0.5, & (0.5)^2 &= 0.25\\
  4.0 - 2.5 &= \phantom{-}1.5, & (1.5)^2 &= 2.25
\end{align*}
\emph{Variance.} Sum $2.25 + 0.25 + 0.25 + 2.25 = 5.00$; divide by 4: $\sigma^2 = 1.25$.

\emph{Standard deviation.} $\sqrt{1.25} = 1.1180$. (With $\epsilon$:
$\sqrt{1.25 + 0.00001} = 1.118038$ against $1.118034$; the difference is in the sixth digit.)

\emph{Normalise.}
\begin{align*}
  -1.5/1.1180 &= -1.3416\\
  -0.5/1.1180 &= -0.4472\\
  \phantom{-}0.5/1.1180 &= \phantom{-}0.4472\\
  \phantom{-}1.5/1.1180 &= \phantom{-}1.3416
\end{align*}
So $\LN(\vect{x}) = (-1.3416, -0.4472, 0.4472, 1.3416)$, which is what PyTorch's
\code{nn.LayerNorm(4)} returns\tagC. The pattern (evenly spaced, increasing) is kept; the level
and the spread are standardised. Check: the outputs sum to 0, and their squares,
$1.8 + 0.2 + 0.2 + 1.8 = 4.0$, average to 1.

\emph{Now with an affine map}, $\boldsymbol\gamma = (2, 1, 1, 0.5)$ and
$\boldsymbol\beta = (0, 0, 1, 0)$:
\begin{align*}
  2\times(-1.3416) + 0 &= -2.6833\\
  1\times(-0.4472) + 0 &= -0.4472\\
  1\times0.4472 + 1 &= \phantom{-}1.4472\\
  0.5\times1.3416 + 0 &= \phantom{-}0.6708
\end{align*}
(using the unrounded $-1.34164$ for the first line). Training can stretch, shrink and shift each
feature after the standardisation.
\end{worked}

\subsection{What LayerNorm forgets}

\begin{keypoint}
LayerNorm is blind to the overall level and overall scale of a token: $\LN(a\vect{x} + b) =
\LN(\vect{x})$ for any number $a > 0$ and any number $b$.
\end{keypoint}

\emph{Proof} (ignoring the tiny $\epsilon$). Let $\vect{y} = a\vect{x} + b$, meaning
$y_i = a x_i + b$ for every $i$.
\begin{align*}
  \mu_y &= \frac{1}{C}\sum_i (a x_i + b) = a\,\frac{1}{C}\sum_i x_i + b = a\mu_x + b && \text{the mean follows the map}\\
  y_i - \mu_y &= a x_i + b - a\mu_x - b = a(x_i - \mu_x) && \text{the shift } b \text{ cancels}\\
  \sigma_y^2 &= \frac{1}{C}\sum_i a^2(x_i - \mu_x)^2 = a^2\sigma_x^2 && \text{square the line above}\\
  \hat y_i &= \frac{a(x_i - \mu_x)}{\sqrt{a^2\sigma_x^2}} = \frac{a(x_i - \mu_x)}{a\,\sigma_x} = \hat x_i && \sqrt{a^2} = a \text{ since } a > 0. \qquad\blacksquare
\end{align*}
The script confirms it: for the token of \cref{we:block:ln},
\[
  \LN(3\vect{x} + 10) = (-1.3416, -0.4472, 0.4472, 1.3416) = \LN(\vect{x})\tagC .
\] Whatever a previous block did to the level and scale
of a token, the next sublayer sees a standardised copy. (The same proof with $C = 1$ shows that a
single number is always mapped to $\beta$; \cref{ex:esa:5}.)

\subsection{LayerNorm against BatchNorm}

\Cref{ch:compressor} met BatchNorm\index{BatchNorm}, which normalises each channel using
statistics gathered over the whole batch. The difference is which numbers are averaged.

\begin{center}
\begin{tabular}{p{0.34\linewidth}p{0.27\linewidth}p{0.27\linewidth}}
\toprule
 & LayerNorm & BatchNorm\\
\midrule
averages over & one token's $C$ features & one channel, across the batch and positions\\
depends on other patches in the batch? & no & yes\\
same in training and evaluation? & yes & no (running statistics)\\
suits token sequences? & yes & poorly\\
\bottomrule
\end{tabular}
\end{center}
With BatchNorm a token's output would depend on which other patches happened to share its batch,
and small batches would give noisy statistics. LayerNorm keeps every token self-contained.
(\Cref{ch:finetune} tells what BatchNorm's batch dependence cost this project in fine-tuning.)

\section{Residual connections}\label{sec:block:residual}

\subsection{The chain rule, on a tiny example}

Training needs the derivative of the loss with respect to every parameter (\cref{ch:trainloop}).
When functions are composed, $y = g(f(x))$, the derivative is the \emph{chain
rule}\index{chain rule}: multiply the derivatives along the way,
\begin{equation}\label{eq:block:chain}
  \frac{dy}{dx} = g'(f(x))\cdot f'(x).
\end{equation}
\emph{Check it on numbers.} Let $f(x) = 3x$ and $g(u) = u^2$, so $y = (3x)^2 = 9x^2$ and directly
$dy/dx = 18x$, which is 36 at $x = 2$. By the chain rule: $f(2) = 6$, $g'(u) = 2u$ so
$g'(6) = 12$, $f'(x) = 3$; then $12\times3 = 36$. The same.

For a stack of $L$ layers $x_\ell = f_\ell(x_{\ell-1})$, the chain rule multiplies $L$ factors:
\begin{equation}\label{eq:block:plainchain}
  \frac{d x_L}{d x_0} = f_L'\cdot f_{L-1}'\cdots f_1' .
\end{equation}

\subsection{What the residual changes}

A \emph{residual connection}\index{residual connection} (or skip connection) computes
\begin{equation}\label{eq:block:res}
  x_\ell = x_{\ell-1} + f_\ell(x_{\ell-1})
\end{equation}
instead of $x_\ell = f_\ell(x_{\ell-1})$ \citep{he2016deep}. Differentiate: the derivative of
$x$ with respect to itself is 1, and the derivative of $f_\ell(x)$ is $f_\ell'$, so
\begin{equation}\label{eq:block:resgrad}
  \frac{d x_\ell}{d x_{\ell-1}} = 1 + f_\ell' .
\end{equation}
Chained over $L$ layers:
\begin{equation}\label{eq:block:reschain}
  \frac{d x_L}{d x_0} = (1 + f_L')(1 + f_{L-1}')\cdots(1 + f_1') .
\end{equation}
Multiply out the brackets and the first term is $1\cdot1\cdots1 = 1$: there is always a path
along which the signal passes untouched, whatever the layers do. (For vectors the 1 becomes the
identity matrix $\mat{I}$ and $f'$ becomes the matrix of partial derivatives, the
\emph{Jacobian}; the conclusion is the same.)

\begin{worked}[label={we:block:resnumbers}]{Twenty layers with and without residuals}
Suppose every layer's derivative is the same number and there are 20 layers.

\emph{Without residuals}, \cref{eq:block:plainchain} gives $(f')^{20}$:
\begin{align*}
  f' = 0.9:&\quad 0.9^{20} = 0.1216\\
  f' = 0.5:&\quad 0.5^{20} = 0.000\,000\,954 .
\end{align*}
Even a mild 0.9 per layer loses 88\% of the signal; 0.5 per layer loses all but one part in a
million. This is the \emph{vanishing gradient}\index{vanishing gradient}.

\emph{With residuals}, \cref{eq:block:reschain} gives $(1 + f')^{20}$. Sublayers start near zero
output, so their derivatives start small:
\begin{align*}
  f' = +0.05:&\quad 1.05^{20} = 2.6533\\
  f' = -0.05:&\quad 0.95^{20} = 0.3585 .
\end{align*}
Both are of order 1. The signal reaching the first layer is neither lost nor exploded.
\end{worked}

\paragraph{The ``learn a correction'' view.} Unroll \cref{eq:block:res} from the input:
\begin{align*}
  x_1 &= x_0 + f_1(x_0)\\
  x_2 &= x_1 + f_2(x_1) = x_0 + f_1(x_0) + f_2(x_1)\\
  &\ \,\vdots\\
  x_L &= x_0 + \sum_{\ell=1}^{L} f_\ell(x_{\ell-1}) .
\end{align*}
The output is the input plus a sum of corrections. A sublayer that has nothing useful to add can
learn $f_\ell = 0$ (all weights near zero), which is easy. Without the residual, ``do nothing''
would mean learning to be exactly the identity function, which is surprisingly hard for a stack of
matrices and nonlinearities.

\section{Pre-norm, not post-norm}\label{sec:block:prenorm}

Where the LayerNorm sits matters.

\begin{center}
\begin{tabular}{lll}
\toprule
 & post-norm (the original 2017 transformer) & pre-norm (this model)\\
\midrule
formula & $x \leftarrow \LN(x + f(x))$ & $x \leftarrow x + f(\LN(x))$\\
the residual path & passes through a LayerNorm every sublayer & untouched\\
\bottomrule
\end{tabular}
\end{center}
\index{pre-norm}\index{post-norm}
In post-norm the derivative of each step is $\LN'\cdot(1 + f')$: the clean ``1'' of
\cref{eq:block:resgrad} is multiplied by LayerNorm's own derivative at every one of the 16
sublayers, so the motorway passes through a toll booth every block. In pre-norm the step is
$x + f(\LN(x))$, whose derivative is $1 + f'\cdot\LN'$: the 1 survives, and the unrolled sum above
holds exactly. Modern transformers (GPT-2, ViT, SegFormer and this code) use pre-norm. The source
says so in its docstring (\file{segformer_block.py}, lines 4--13).

\section{Dropout}\label{sec:block:dropout}

\subsection{What it does}

\emph{Dropout}\index{dropout} \citep{srivastava2014dropout}, during training, sets each number of a
tensor to zero independently with probability $p$, and multiplies the survivors by
$1/(1-p)$. At evaluation it does nothing.

\emph{Why the factor $1/(1-p)$.} Write the dropped value as $m\,x/(1-p)$ where $m$ is 1 with
probability $1-p$ and 0 with probability $p$. Its expected value is
\begin{align*}
  \E\!\left[\frac{m\,x}{1-p}\right] = \frac{x}{1-p}\,\E[m] = \frac{x}{1-p}\,\big((1-p)\cdot1 + p\cdot0\big) = \frac{x}{1-p}(1-p) = x .
\end{align*}
On average the output equals the input, so switching dropout off at evaluation changes nothing on
average.

\begin{worked}[label={we:block:dropout}]{Dropout with $p = 0.4$}
Scale factor $1/(1-0.4) = 1/0.6 = 1.6667$. Take $x = (0.5, -1.0, 2.0, 0.3, 1.2)$ and suppose the
random draw keeps positions 1, 3 and 4 ($m = (1, 0, 1, 1, 0)$):
\begin{align*}
  0.5\times1\times1.6667 &= 0.8333\\
  -1.0\times0 &= 0\\
  2.0\times1\times1.6667 &= 3.3333\\
  0.3\times1\times1.6667 &= 0.5000\\
  1.2\times0 &= 0 .
\end{align*}
Two of five values vanished; the survivors grew by two thirds.
\end{worked}

\subsection{Where it sits}

The block applies dropout to the sublayer's output \emph{before} the residual addition,
\[
  x + \operatorname{Dropout}\big(f(\LN(x))\big).
\]
Only the new contribution can be dropped; the information
travelling along the skip connection is never touched. Dropping the residual stream itself would
destroy accumulated information at random, which is damage, not regularisation.

\begin{drift}
Dropout rates (\code{drop_rate}): 0.3 in v0.1.0
\tagM{configs/hyperspectral_segformer_exp_1.json} and 0.4 from v0.2.0 on
\tagM{configs/hyperspectral_segformer_exp_2.json}, also 0.4 in FT04
\tagM{research/experiments/indradhanu-finetune/FT04-final-joint/config.json}. At 0.4 nearly half of
every sublayer's output is discarded in training; it was raised to fight overfitting. At inference
dropout is off twice over: the model is built with \code{drop_rate=0.0} and put in
\code{model.eval()} mode (IW15).
\end{drift}

\begin{sourcenote}[: one dropout module, and no drop path]
\file{segformer_block.py} builds a single \code{self.drop = nn.Dropout(drop_ratio)} (line 79) and
uses it after both sublayers (lines 96 and 103). That is harmless: dropout has no parameters and
draws a fresh random mask on every call. The block has no \emph{drop path} (stochastic depth, which
drops a whole sublayer's contribution for a whole sample and which many transformer implementations
add): a search of \file{app/} for \code{DropPath} or \code{drop_path} finds nothing. The only
regularising randomness inside the block is this element-wise dropout.
\end{sourcenote}

\section{Mix-FFN}\label{sec:block:mixffn}

\subsection{The feed-forward network, and why it widens}

The standard transformer feed-forward network\index{feed-forward network} processes each token on
its own:
\[
  \operatorname{Linear}(C \to EC) \;\to\; \text{nonlinearity} \;\to\; \operatorname{Linear}(EC \to C),
\]
with an \emph{expansion ratio}\index{expansion ratio} $E = 4$ here
\par\noindent\tagM{research/experiments/indradhanu-finetune/FT04-final-joint/config.json}.
\par\noindent The middle is four times wider than the ends. A wider middle gives the nonlinearity more room: the token is unfolded
into $4C$ numbers, each of which can be bent independently, and then folded back. Going from 32 to
32 through one nonlinearity is far more restrictive than from 32 to 128 and back.

\subsection{SegFormer's twist: a depthwise convolution in the middle}

SegFormer's Mix-FFN\index{Mix-FFN} \citep{xie2021segformer} inserts a $3\times3$ convolution on the
token grid between the expansion and the nonlinearity:
\begin{align*}
  (B, N, C) &\xrightarrow{\ \text{fc1: Linear}(C\to4C)\ } (B, N, 4C)
  \xrightarrow{\ \text{transpose, reshape}\ } (B, 4C, H, W)\\
  &\xrightarrow{\ \text{dwconv } 3\times3\ } (B, 4C, H, W)
  \xrightarrow{\ \text{reshape, transpose}\ } (B, N, 4C)\\
  &\xrightarrow{\ \GELU\ } (B, N, 4C)
  \xrightarrow{\ \text{fc2: Linear}(4C\to C)\ } (B, N, C).
\end{align*}
The convolution has stride 1 and padding 1, so by \cref{eq:esa:outsize} the grid keeps its size:
$\lfloor(H + 2 - 3)/1\rfloor + 1 = H$. \emph{Mix-FFN never changes the number of tokens.}

\coursefig[0.95\linewidth]{iw-14-mixffn}{The feed-forward half of a block. Tokens are expanded
fourfold, reshaped back to their grid, mixed with their eight neighbours by a depthwise
$3\times3$ convolution (one private filter per channel), passed through GELU and projected back.
The depthwise convolution is what gives the model a sense of position.}{fig:block:mixffn}

\subsection{Depthwise convolution}

\begin{definition}[Depthwise convolution]
A convolution with \code{groups} equal to the number of channels: each channel has its own private
$3\times3$ filter and never looks at any other channel,
\begin{equation}\label{eq:block:dw}
  y_{c,i,j} = b_c + \sum_{a=-1}^{1}\sum_{b=-1}^{1} w_{c,a,b}\; x_{c,\,i+a,\,j+b},
\end{equation}
with $x = 0$ outside the grid (zero padding).
\end{definition}\index{depthwise convolution}\index{convolution!depthwise}
In code (\file{mix_ffn.py}, line 79):
\begin{lstlisting}
self.dwconv = nn.Conv2d(hidden_dim, hidden_dim, kernel_size=3,
                        padding=1, groups=hidden_dim)   # (B, 4C, H, W) -> (B, 4C, H, W)
\end{lstlisting}
\emph{Parameters.} An ordinary $3\times3$ convolution from $4C$ channels to $4C$ channels has
$4C\cdot4C\cdot9 + 4C$; the depthwise one has 9 weights per channel and a bias:
\begin{equation}\label{eq:block:dwparams}
  \text{dwconv} = 4C\cdot9 + 4C = 40C .
\end{equation}

\begin{worked}[label={we:block:dwcount}]{Ordinary against depthwise at stage 1}
Stage 1 has $C = 32$, so $4C = 128$ hidden channels.
\begin{align*}
  \text{ordinary: } 128\times128\times9 + 128 &= 16{,}384\times9 + 128 = 147{,}456 + 128 = 147{,}584\\
  \text{depthwise: } 128\times9 + 128 &= 1{,}152 + 128 = 1{,}280 .
\end{align*}
$147{,}584/1{,}280 = 115$ times fewer. Nothing is lost, because \code{fc1} and \code{fc2} on either
side already mix channels; the convolution's job is purely spatial.
\end{worked}

\subsection{Where position comes from}

\Cref{sec:attention:perm} showed that attention cannot tell where a token is. A convolution can:
its output at $(i, j)$ depends on the tokens at $(i\pm1, j\pm1)$, so ``next to'' is built into it.
More than that, the zero padding makes the border different from the interior, and over several
layers that difference spreads inward, so the network can tell how far a token is from the edge.
SegFormer's authors argue that this $3\times3$ convolution provides enough positional information
to do without positional encodings \citep{xie2021segformer}; this codebase has none.

\begin{worked}[label={we:block:dwtoy}]{A depthwise filter sees the border}
One channel on a $3\times3$ grid, values rising evenly:
\[
  \begin{pmatrix} 1 & 2 & 3\\ 4 & 5 & 6\\ 7 & 8 & 9 \end{pmatrix},
  \qquad \text{filter } \begin{pmatrix} 0 & 1 & 0\\ 1 & -4 & 1\\ 0 & 1 & 0 \end{pmatrix},\ \text{bias } 0 .
\]
The filter adds the four direct neighbours and subtracts four times the centre.

\emph{Centre $(2,2)$.} Neighbours 2 (above), 4 (left), 6 (right), 8 (below):
\[
  2 + 4 + 6 + 8 - 4\times5 = 20 - 20 = 0 .
\]
\emph{Top-left corner $(1,1)$.} Above and left are padding zeros; right is 2, below is 4:
\[
  0 + 0 + 2 + 4 - 4\times1 = 6 - 4 = 2 .
\]
\emph{Top-middle $(1,2)$.} Above is padding; left 1, right 3, below 5:
\[
  0 + 1 + 3 + 5 - 4\times2 = 9 - 8 = 1 .
\]
The whole output, from the script\tagC:
\[
  \begin{pmatrix} 2 & 1 & -4\\ -3 & 0 & -7\\ -16 & -11 & -22 \end{pmatrix}.
\]
On a perfectly even ramp the interior gives 0, but every border position gives something else,
and each corner something different. The output carries a trace of \emph{where} each token is,
which no amount of attention could produce.
\end{worked}

\subsection{GELU}

The nonlinearity is the Gaussian Error Linear Unit\index{GELU} \citep{hendrycks2016gaussian}:
\begin{equation}\label{eq:block:gelu}
  \GELU(x) = x\,\Phi(x),\qquad \Phi(x) = \tfrac12\left(1 + \operatorname{erf}\!\left(\tfrac{x}{\sqrt2}\right)\right),
\end{equation}
where $\Phi(x)$ is the probability that a standard Gaussian random number is below $x$ (the
Gaussian's cumulative distribution function; $\operatorname{erf}$ is the error function that
computes it). Read it as ``let $x$ through in proportion to how likely a Gaussian is to be below
it'': large positive $x$ passes almost whole ($\Phi \approx 1$), large negative $x$ is almost
blocked ($\Phi\approx 0$), and in between it bends smoothly. PyTorch's \code{nn.GELU()} (used at
\file{mix_ffn.py}, line 86) computes this exact form; a common approximation is
$\tfrac12 x\big(1 + \tanh(\sqrt{2/\pi}\,(x + 0.044715x^3))\big)$.

\begin{worked}[label={we:block:gelu}]{GELU at $-1$, 0 and 1}
From tables of the Gaussian, $\Phi(1) = 0.8413$, and by symmetry $\Phi(-1) = 1 - 0.8413 =
0.1587$; $\Phi(0) = 0.5$.
\begin{align*}
  \GELU(-1) &= -1\times0.1587 = -0.1587\\
  \GELU(0) &= 0\times0.5 = 0\\
  \GELU(1) &= 1\times0.8413 = 0.8413 .
\end{align*}
PyTorch gives $(-0.1587, 0, 0.8413)$\tagC; the tanh approximation gives
$(-0.1588, 0, 0.8412)$.

\emph{The slope.} By the product rule, $\GELU'(x) = \Phi(x) + x\,\varphi(x)$, where
$\varphi(x) = e^{-x^2/2}/\sqrt{2\pi}$ is the Gaussian density and $\varphi(1) = \varphi(-1) =
0.2420$:
\begin{align*}
  \GELU'(-1) &= 0.1587 + (-1)(0.2420) = -0.0833\\
  \GELU'(0) &= 0.5 + 0 = 0.5\\
  \GELU'(1) &= 0.8413 + 0.2420 = 1.0833 .
\end{align*}
Unlike ReLU, whose slope jumps from 0 to 1 at the origin, GELU's slope changes smoothly; it even
dips slightly below zero for negative inputs, so GELU is not monotonic near $x = -0.75$.
\end{worked}

\begin{figure}[htbp]
\centering
\begin{tikzpicture}
\begin{axis}[width=0.55\linewidth, height=6cm, xmin=-3, xmax=3, ymin=-0.6, ymax=3,
  axis lines=middle, xlabel={$x$}, ylabel={output},
  legend style={font=\small, at={(1.04,0.5)}, anchor=west}, grid=major, grid style={gray!20}]
\addplot[maskgrey, very thick, dashed, domain=-3:3, samples=61] {max(0,x)};
\addplot[spat, very thick, domain=-3:3, samples=121] {0.5*x*(1+tanh(0.7978845608*(x+0.044715*x^3)))};
\addplot[only marks, mark=*, attn] coordinates {(-1,-0.1587) (0,0) (1,0.8413)};
\legend{{ReLU $\max(0,x)$}, GELU, {values of \cref{we:block:gelu}}}
\end{axis}
\end{tikzpicture}
\caption{GELU against ReLU. GELU is a smooth ReLU: near zero for very negative inputs, a soft bend
through the origin, near-linear for large inputs. The curve is drawn with the tanh approximation,
which differs from the exact form by less than 0.001 on this range.}
\label{fig:block:gelu}
\end{figure}

GELU sits \emph{after} the depthwise convolution and \emph{before} \code{fc2}, that is, in the wide
$4C$-dimensional space, where the nonlinearity has the most numbers to bend.

\subsection{The sparse fallback}

Like the reduction of \cref{sec:esa:fallback}, the depthwise convolution needs a full grid,
$N = H\cdot W$. When stage-1 tokens have been removed it is skipped (\file{mix_ffn.py}, line 119):
\begin{lstlisting}
if N == H * W:
    ...                       # reshape to grid, dwconv, reshape back: (B, N, 4C)
else:
    sequence_x = expanded_x   # (B, N, 4C): no spatial mixing
\end{lstlisting}
In a masked pass, stage 1's Mix-FFN is a plain two-layer network with a GELU in the middle. The
docstring calls this ``a minor quality tradeoff for correct handling of sparse token sets'', since
the attention still supplies spatial context. As with the reduction convolution, the stage-1
depthwise filters receive no gradient from a masked pass.

\subsection{Mix-FFN parameter count}

Three layers: \code{fc1} has $C\cdot4C + 4C$, \code{dwconv} has $40C$ (\cref{eq:block:dwparams}),
\code{fc2} has $4C\cdot C + C$. In total
\begin{equation}\label{eq:block:ffnparams}
  \text{Mix-FFN}(C) = (4C^2 + 4C) + 40C + (4C^2 + C) = 8C^2 + 45C .
\end{equation}

\begin{worked}[label={we:block:ffn1}]{Mix-FFN at stage 1}
$C = 32$, hidden width $4\times32 = 128$.
\begin{align*}
  \text{fc1: } 32\times128 + 128 &= 4{,}096 + 128 = 4{,}224\\
  \text{dwconv: } 128\times9 + 128 &= 1{,}152 + 128 = 1{,}280\\
  \text{fc2: } 128\times32 + 32 &= 4{,}096 + 32 = 4{,}128\\
  \text{total: } 4{,}224 + 1{,}280 + 4{,}128 &= 9{,}632 .
\end{align*}
Check with \cref{eq:block:ffnparams}: $8\times32^2 + 45\times32 = 8\times1{,}024 + 1{,}440 =
8{,}192 + 1{,}440 = 9{,}632$.
\end{worked}

\section{The block}\label{sec:block:block}

\subsection{Two lines}

Putting the pieces together (\file{segformer_block.py}, lines 95--104):
\begin{align}
  \vect{a} &= \vect{x} + \operatorname{Dropout}\big(\operatorname{ESA}(\LN_1(\vect{x}), H, W)\big), \label{eq:block:sub1}\\
  \vect{y} &= \vect{a} + \operatorname{Dropout}\big(\operatorname{MixFFN}(\LN_2(\vect{a}), H, W)\big). \label{eq:block:sub2}
\end{align}
\index{SegFormer block}\index{transformer block}
\begin{lstlisting}
attn = self.drop(self.attn(self.norm1(x), H, W))   # (B, N, C)
attn = x + attn                                    # (B, N, C) residual 1
ffn  = self.drop(self.ffn(self.norm2(attn), H, W)) # (B, N, C)
ffn  = attn + ffn                                  # (B, N, C) residual 2
\end{lstlisting}
Read it aloud: \emph{normalise, attend, maybe drop some, add it to what we had; then normalise,
think, maybe drop some, add it to what we had.} The second residual adds to $\vect{a}$, the output
of the first sublayer, not to the original $\vect{x}$; information accumulates through both. The
shape \shape{(B, N, C)} never changes inside a block.

\coursefig[0.95\linewidth]{iw-15-block}{One SegFormer block. Each of the two sublayers, efficient
self-attention (plum) and Mix-FFN (indigo), normalises its input first (pre-norm), transforms it,
applies dropout and adds the result back to its input through a residual connection. The shape
$(B, N, C)$ is unchanged; each encoder stage stacks two blocks.}{fig:block:course}

Each stage stacks \code{num_blocks} $= [2, 2, 2, 2]$ blocks, the SegFormer-B0 setting
\par\noindent\tagM{research/experiments/indradhanu-finetune/FT04-final-joint/config.json}.
\par\noindent Within a stage all blocks share $C$, $h$ and $R$; width and resolution change only at the
patch embeddings between stages (\cref{ch:encoder}).

\subsection{The parameter count, by hand}\label{sec:block:params}

A block is two LayerNorms, one ESA and one Mix-FFN. With
\cref{eq:esa:params,eq:block:lnparams,eq:block:ffnparams}:
\begin{equation}\label{eq:block:blockparams}
  \text{block}(C, R) = \underbrace{2\cdot2C}_{\text{norm1, norm2}} + \underbrace{4C^2 + 4C + [R > 1](C^2R^2 + 3C)}_{\text{ESA}} + \underbrace{8C^2 + 45C}_{\text{Mix-FFN}},
\end{equation}
where $[R > 1]$ is 1 if $R > 1$ and 0 otherwise.

\begin{worked}[label={we:block:stage1}]{A stage-1 block: 79{,}616 parameters}
$C = 32$, $R = 8$, $h = 2$ (the head count does not matter, \cref{sec:attention:heads}).
\begin{align*}
  \text{norm1: } 2\times32 &= 64\\
  \text{q, k, v, proj: } 4\times(32\times32 + 32) = 4\times1{,}056 &= 4{,}224\\
  \text{reduction conv: } 32\times32\times8\times8 + 32 = 65{,}536 + 32 &= 65{,}568\\
  \text{reduction norm: } 2\times32 &= 64\\
  \text{norm2: } 2\times32 &= 64\\
  \text{fc1: } 32\times128 + 128 &= 4{,}224\\
  \text{dwconv: } 128\times9 + 128 &= 1{,}280\\
  \text{fc2: } 128\times32 + 32 &= 4{,}128
\end{align*}
Add, in groups: ESA $= 4{,}224 + 65{,}568 + 64 = 69{,}856$; Mix-FFN $= 4{,}224 + 1{,}280 + 4{,}128 =
9{,}632$; the two block norms $= 64 + 64 = 128$. One block:
\[
  69{,}856 + 9{,}632 + 128 = 79{,}616 .
\]
Two blocks per stage: $2\times79{,}616 = 159{,}232$.

\emph{Check against the real model.} The \code{torchinfo} summary of the hyperspectral SegFormer
MAE lists stage 1's block list (\code{ModuleList: 3-2}) at 159{,}232 parameters
\tagM{research/model_break_down/05_hyperspectral_segformer_mae.md}, and the course on Indradhanu
Heavy quotes FT04's stage-1 block as 79{,}616 \tagM{indradhanu workings/68-indradhanu-heavy-pretraining.md}.
Building \code{SegFormerBlock(32, 2, 8)} in the script gives 79{,}616 with exactly the
eight per-layer counts above\tagC. You have predicted a number from a real model by adding up
pieces derived by hand in three chapters.
\end{worked}

\begin{worked}[label={we:block:stage4}]{A stage-4 block: 800{,}000 parameters}
$C = 256$, $R = 1$ (no reduction layers), hidden width $4\times256 = 1{,}024$.
\begin{align*}
  \text{norm1: } 2\times256 &= 512\\
  \text{q, k, v, proj: } 4\times(256\times256 + 256) = 4\times(65{,}536 + 256) = 4\times65{,}792 &= 263{,}168\\
  \text{norm2: } 2\times256 &= 512\\
  \text{fc1: } 256\times1{,}024 + 1{,}024 = 262{,}144 + 1{,}024 &= 263{,}168\\
  \text{dwconv: } 1{,}024\times9 + 1{,}024 = 9{,}216 + 1{,}024 &= 10{,}240\\
  \text{fc2: } 1{,}024\times256 + 256 = 262{,}144 + 256 &= 262{,}400
\end{align*}
ESA $= 263{,}168$; Mix-FFN $= 263{,}168 + 10{,}240 + 262{,}400 = 535{,}808$; norms $= 512 + 512 =
1{,}024$. One block:
\[
  263{,}168 + 535{,}808 + 1{,}024 = 800{,}000 .
\]
A round number by coincidence. Two blocks: 1{,}600{,}000, which is the \code{torchinfo} figure for
stage 4's block list \tagM{research/model_break_down/05_hyperspectral_segformer_mae.md}, and the
script's\tagC.
\end{worked}

All four stages, checked against the built blocks\tagC\ and the \code{torchinfo} table
\par\noindent\tagM{research/model_break_down/05_hyperspectral_segformer_mae.md}:
\begin{center}
\begin{tabular}{crrrrrr}
\toprule
stage & $C$ & ESA & Mix-FFN & 2 LNs & one block & two blocks\\
\midrule
1 & 32 & 69{,}856 & 9{,}632 & 128 & 79{,}616 & 159{,}232\\
2 & 64 & 82{,}368 & 35{,}648 & 256 & 118{,}272 & 236{,}544\\
3 & 160 & 205{,}920 & 212{,}000 & 640 & 418{,}560 & 837{,}120\\
4 & 256 & 263{,}168 & 535{,}808 & 1{,}024 & 800{,}000 & 1{,}600{,}000\\
\midrule
\multicolumn{6}{l}{all eight blocks} & 2{,}832{,}896\\
\bottomrule
\end{tabular}
\end{center}

\paragraph{The balance flips.} At stage 1 attention is 88\% of the block ($69{,}856/79{,}616$),
almost all of it the $8\times8$ reduction convolution. At stage 4 Mix-FFN is 67\%
($535{,}808/800{,}000$): with no reduction and $C = 256$, the two $C\times4C$ matrices dominate.
\Cref{ch:params} puts these blocks into the budget of the whole model.

\begin{drift}
Indradhanu Heavy (\cref{ch:heavy}) uses the same block with $C = 96$ at stage 1; its stage-1 block
is 705{,}792 parameters \tagM{indradhanu workings/68-indradhanu-heavy-pretraining.md}, which the
script reproduces\tagC. You count it in \cref{ex:block:5}.
\end{drift}

\begin{confusions}
\confusion{Why two LayerNorms per block?}{One before attention, one before Mix-FFN. Each sublayer
gets a freshly standardised copy of its input. The ESA's reduction has a third LayerNorm of its own,
on the summary tokens.}
\confusion{Is the plus sign ordinary addition?}{Yes: element-by-element addition of two tensors
of the same shape \texttt{(B, N, C)}. Its power is in the derivative, $1 + f'$.}
\confusion{Is Mix-FFN mixing tokens or channels?}{Both, in different places. \texttt{fc1} and
\texttt{fc2} mix channels within each token; the depthwise convolution mixes neighbouring tokens
within each channel. Between them everything is mixed, hence the name.}
\confusion{Does Mix-FFN change the number of tokens?}{No. Stride 1 and padding 1 keep the grid
size. Only the patch embeddings change the token count.}
\confusion{LayerNorm or BatchNorm: which normalises over the batch?}{BatchNorm. LayerNorm uses one
token's own $C$ features and nothing else, and so behaves the same in training and evaluation.}
\confusion{Does dropout run at inference?}{No. \texttt{drop\_rate=0.0} and \texttt{model.eval()}
both switch it off. Run-to-run variation at inference comes from the masking strategy
(\cref{ch:inference}), not from dropout.}
\confusion{Is GELU the same as ReLU?}{It is a smooth relative of ReLU. It equals ReLU far from zero, but
it is curved near zero, its slope changes continuously, and it lets small negative values through
(GELU$(-1) = -0.1587$).}
\end{confusions}

\begin{checkyourself}
\item\label{ex:block:1} LayerNorm the token $(2.0, 4.0, 4.0, 6.0)$ with $\gamma = 1$, $\beta = 0$
(ignore $\epsilon$). Show the mean, every deviation and square, the variance, the standard deviation
and the four outputs. Then predict, without calculating, $\LN(10, 20, 20, 30)$.
\item\label{ex:block:2} Compute $\GELU(2)$ and $\GELU(-2)$ given $\Phi(2) = 0.9772$. What fraction
of the input does GELU pass at $x = 2$ and at $x = -2$?
\item\label{ex:block:3} A stack of 16 sublayers (the encoder's eight blocks) each has derivative
$f' = -0.1$. Compute the end-to-end derivative without residuals and with pre-norm residuals
(assume LayerNorm's derivative is absorbed into $f'$). Which one still lets the first layer learn?
\item\label{ex:block:4} Count the Mix-FFN parameters for stage 2 ($C = 64$) layer by layer, and
check your answer with $8C^2 + 45C$.
\item\label{ex:block:5} Indradhanu Heavy's stage-1 block has $C = 96$, $R = 8$, expansion 4. Count
its parameters layer by layer and confirm 705{,}792. Why is it close to nine times FT04's 79{,}616
but not exactly nine?
\item\label{ex:block:6} With dropout probability $p = 0.4$, what is the expected output and the
variance of the dropped value $m\,x/(1-p)$ for an input $x = 1$? (Use
$\operatorname{Var} = \E[\cdot^2] - (\E[\cdot])^2$.)
\end{checkyourself}

\nextchapter{ch:encoder}
